\section{Conclusions}\label{sec:conclusions}

We have presented the universal Wilson coefficients for the
gradient-flow definition of the energy-momentum tensor through
\nnlo\ \qcd. The \nnlo\ corrections modify the three numerically
dominant coefficients $c_1$, $c_2$, $\ringc_3$ at the level of 10\%
(1-2\%) for a central scale of $\mu_0=3$\,GeV ($\mu_0=130$\,GeV), where
$\mu_0$ is related to the flow time $t$ according to \eqn{eq:mu0}. We
observe a reduction of the theoretical uncertainty relative to the
\nlo\ result as derived from varying the renormalization scale by a
factor of two around its central value.  The behavior of the fourth
coefficient $\ringc_4$ is less satisfactory, but its impact is expected
to be numerically suppressed.

Aside from this main outcome, new results presented in this paper
include the flowed quark-field renormalization constant to \nnlo\ in the
\msbar\ scheme, and the anomalous dimension matrix for the regular
\qcd\ operators which make up the \emt.

In conclusion, we hope that our results will help to improve the studies
of the \emt\ on the lattice. They are the first outcome of a systematic
setup for higher-order perturbative calculations within the
gradient-flow formalism\,\cite{Artz:2019bpr}, which should prove useful also
in other applications of this theoretical framework.
